\documentclass[12pt]{article}
\usepackage[utf8]{inputenc}
\usepackage[T1]{fontenc}
\usepackage{hyperref}
\title{Thachanadan Language Resources for K-12 Curriculum Development}
\author{Compiled by Claude Sonnet 5 (claude-sonnet-5) for Andr\'as Kornai}
\date{}
\begin{document}
\maketitle

\section{Introduction}
Thachanadan (ISO 639-3 thn) is spoken by the Thachanadan Moopan (Muppan) community. The sources seen place it in Wayanad district, Kerala, with about 1,649 people counted in 2008, and Wikipedia reports about 3,000 speakers in a 2004 survey. It is classified as Southern Dravidian, though SPPEL describes it as a mix of Malayalam, Telugu and Tulu. It has no script and is mostly spoken by elders, as younger people shift to Malayalam. Very little material exists.

\section{Primary Educational Resources}
None found.

\section{Linguistic and Reference Materials}
\begin{itemize}
\item Wikipedia entry on Thachanadan: \url{https://en.wikipedia.org/wiki/Thachanadan_language}
\item CIIL Sanchika special collection: \url{https://sanchika.ciil.org/communities/d4941508-a0b2-4b08-8f4f-8f694b3fbc94}
\item CIIL Sanchika record: \url{https://sanchika.ciil.org/handle/1/345}
\end{itemize}

\section{Religious and Community Resources}
\begin{itemize}
\item SPPEL community documentation on Thachanadan Moopan: \url{https://www.sppel.org/thachanadanmoopandoc.aspx}
\item Joshua Project language profile: \url{https://joshuaproject.net/languages/thn}
\end{itemize}

\section{Digital and Audio Resources}
\begin{itemize}
\item SPPEL language documentation page: \url{https://sppel.org/resources/language-documentation/thachanadan-moopan}
\end{itemize}
The contents of the SPPEL and CIIL pages were not inspected beyond their listing, so they may hold recordings or wordlists.

\section{Curriculum Development Context}
There is no known written tradition or school programme. A curriculum would have to begin with documentation and a community-agreed orthography, probably using Malayalam script, while Malayalam serves as the school language.
\end{document}